\begin{textsample}{POS Dim 7 – global\_south\_2024\_09 – Score 1.00 – global\_south\_2024\_09/114426754.txt}  \label{ex:f7_pos_279}
\textbf{Saudi} Crown Prince says no to Israel ties without Palestinian state"The kingdom will not cease its tireless efforts to establish an independent Palestinian state with East Jerusalem as its capital, and we affirm that the kingdom will not establish diplomatic relations with Israel without one,"the crown prince and the kingdom ' s de facto ruler said Wednesday in an address to a senior advisory council.
Crown Prince Mohammed bin Salman of \textbf{Saudi} Arabia has declared that the kingdom will not establish diplomatic relations with Israel before the"establishment of a Palestinian state,"an apparent hardening of his position on an issue that could reshape the diplomatic map of the Middle East."The kingdom will not cease its tireless efforts to establish an independent Palestinian state with East Jerusalem as its capital, and we affirm that the kingdom will not establish diplomatic relations with Israel without one,"the crown prince and the kingdom ' s de facto ruler said Wednesday in an address to a senior advisory council."We thank all the countries that recognized the Palestinian state as an embodiment @ @ @ @ @ @ @ @ @ @ take similar steps."For decades, the leaders of \textbf{Saudi} Arabia, like those of most other Arab countries, refused to recognize Israel without the creation of a state for the Palestinians.
But after 2020, when four Arab states established formal ties with Israel in agreements brokered by then-President Donald Trump, Crown Prince Mohammed became the first \textbf{Saudi} leader to talk openly about the possibility of \textbf{Saudi} Arabia doing the same.
Advertisement In an interview with Fox News last September, he called a potential agreement"the biggest historical deal since the end of the Cold War"and said it would require"a good life for the Palestinians"but did not mention Palestinian statehood.
His statement Wednesday followed a general hardening of official \textbf{Saudi} rhetoric toward Israel since the start of the Israel-Hamas war in October."We renew the kingdom ' s rejection and strong condemnation of the crimes of the Israeli occupation authority against the Palestinian people,"Crown Prince Mohammed said, delivering remarks on behalf of his father @ @ @ @ @ @ @ @ @ @ month after President Mahmoud Abbas of the Palestinian Authority visited him in Riyadh, the capital of \textbf{Saudi} Arabia.
Advertisement Until Hamas sparked the war in the Gaza Strip with a devastating attack on Israel on Oct. 7, both Israeli and \textbf{Saudi} officials had been indicating they were moving toward a deal. \textbf{Saudi} Arabia has been seeking security guarantees, including a defense pact with the United States and assistance with a civilian nuclear program as part of any agreement.
During the Fox News interview, Crown Prince Mohammed stated that \textbf{Saudi} Arabia was"getting closer"to an accord.
However, since the war in Gaza broke out, \textbf{Saudi} Arabia has insisted on the need for an"irreversible track"to Palestinian statehood.
This month, Secretary of State Antony Blinken said he still hoped to finalize an Israeli-Saudi deal before the end of President Joe Biden ' s term."I think if we can get a cease-fire in Gaza, there remains an opportunity through the balance of this administration to move forward @ @ @ @ @ @ @ @ @ @ Palestinian officials welcomed the crown prince ' s comments, saying they supported the position of the Ramallah, West Bank-based Palestinian leadership."This is an affirmation that the \textbf{Saudi} position is enduring in its support of the Palestinian cause,"Mahmoud al-Habbash, the religious affairs adviser to Abbas, said in an interview Thursday."We ' re reassured about the \textbf{Saudi} stance, which is a cornerstone of the Arab and Islamic world ' s position."

% matched lemmas: saudi
\end{textsample}
